\documentclass[11pt]{article}

\usepackage[a4paper,margin=0.78in]{geometry}
\usepackage[T1]{fontenc}
\usepackage[utf8]{inputenc}
\usepackage{lmodern}
\usepackage{amsmath,amssymb,amsthm,mathtools,bm}
\usepackage{enumitem,booktabs,array,tabularx,microtype,xcolor}
\usepackage[numbers,sort&compress]{natbib}
\usepackage[colorlinks=true,linkcolor=blue!65!black,
  citecolor=blue!65!black,urlcolor=blue!70!black]{hyperref}
\usepackage[nameinlink,noabbrev]{cleveref}
\setlength{\emergencystretch}{3em}

\hypersetup{
  pdftitle={The Canonical Cut-Occurrence Shadow},
  pdfauthor={Liang Wang},
  pdfsubject={A maximal source-backed partial lift and its exact boundary},
  pdfkeywords={fixed shift, occurrence lift, frontier, provenance, conservation}
}

\newtheorem{theorem}{Theorem}[section]
\newtheorem{proposition}[theorem]{Proposition}
\newtheorem{corollary}[theorem]{Corollary}
\theoremstyle{definition}
\newtheorem{definition}[theorem]{Definition}
\theoremstyle{remark}
\newtheorem{remark}[theorem]{Remark}

\newcommand{\ETO}{\mathsf{ETO}}
\newcommand{\FUM}{\mathsf{FUM}}
\newcommand{\EPS}{\mathsf{EPS}}
\newcommand{\Cns}{\mathcal C_X^{\rm ns}}
\newcommand{\Ocut}{\mathcal O_X^{\rm cut}}
\newcommand{\Oact}{\mathcal O_X^{\rm act}}
\newcommand{\one}{\mathbf 1}
\newcommand{\NT}{\textnormal{\textsc{not-testable}}}
\newcommand{\PROVED}{\textnormal{\textsc{proved}}}
\newcommand{\STOP}{\textnormal{\textsc{stop-declared-route}}}
\newcommand{\Lzero}{\mathrm{L0}}
\newcommand{\Lone}{\mathrm{L1}}
\newcommand{\Ltwo}{\mathrm{L2}}
\newcolumntype{Y}{>{\raggedright\arraybackslash}X}
\newcolumntype{P}[1]{>{\raggedright\arraybackslash}p{#1}}

\title{\textbf{The Canonical Cut-Occurrence Shadow:}\\
\textbf{A Maximal Source-Backed Partial Lift}\\
\textbf{and Its Exact Boundary}}
\author{Liang Wang\\
School of Artificial Intelligence and Automation,\\
Huazhong University of Science and Technology, Wuhan 430074, P.R. China\\
\texttt{wangliang.f@gmail.com}}
\date{July 2026}

\begin{document}
\maketitle

\begin{abstract}
The first unsupported object in the source-locked TPC route is a
conservative, row-separated lift from every nonsoft cut path to its
downstream occurrences.  The existing archive does not contain those
occurrences.  We construct instead the strongest object that is
literally determined by the archive: the canonical cut-occurrence
shadow
\[
 S_X:\mathbb C^{\mathcal C_X^{\rm ns}}
 \longrightarrow\mathbb C^{\mathcal O_X^{\rm cut}},
 \qquad S_Xe_c=e_{\iota(c)}.
\]
There is exactly one shadow row for every eligible-tail-open or
frontier-unmapped cut path.  Thus \(S_X\) is injective, row separated
and exactly conservative, \(\mathbf1^TS_X=\mathbf1^T\).  It preserves
the native tuple, prescribed shift, physical normalization, symbolic
coefficient and cut multiplier, while keeping canonical-parent,
stage, physical-occurrence and active-support fields absent.  We
prove a conditional pushforward property: every future actual
occurrence lift that faithfully extends the cut data must forget to
\(S_X\).  This universal property does not produce such a lift.
The deterministic certificate reruns the TPC-143 source chain and
contains \(2988\) production rows, all frontier-unmapped.  A separate
eligible-tail fixture is marked synthetic \(\mathrm{L0}\) only and
never enters the production census.  The result is structural
\(\mathrm{L1}\) progress, not a fixed-shift arithmetic saving.
\end{abstract}

\section{The exact input}

The native, dyadic and cut archives of TPC-133--136 enumerate a
literal packet through the first unsupported carrier
\citep{WangTPC133,WangTPC134,WangTPC135,WangTPC136}.  Their three
terminal types are
\[
 \EPS,\qquad \ETO,\qquad \FUM.
\]
The already soft prefix rows \(\EPS\) need no occurrence repair at
this gate.  The required nonsoft domain is
\begin{equation}
 \Cns=\mathcal C_X^{\ETO}\sqcup\mathcal C_X^{\FUM}.
 \label{eq:nonsoft-domain}
\end{equation}
Neither summand may be omitted merely because it is empty in one
finite fixture.  TPC-143 attaches an explicit lift obligation to
every element of \(\Cns\), but correctly leaves every actual
occurrence edge absent \citep{WangTPC143}.

For \(c\in\Cns\), let \(F_{\rm cut}(c)\) denote precisely the sourced
data available at the cut:
\begin{enumerate}[label=(\roman*),leftmargin=2.2em]
\item the cut, upstream-path and native identities;
\item \(X,h_0,Q,U,V\), the weight-source identifier and physical
normalization;
\item the native tuple \((\ell,k,d_{\rm nat})\);
\item the native symbolic coefficient, native constraint witnesses
and dyadic multiplier expressions; and
\item the terminal type, block, boundary rule and formal-support
status.
\end{enumerate}
The subscript in \(d_{\rm nat}\) is deliberate.  No existing source
identifies this divisor coordinate with a later affine intercept.

\begin{definition}[Support namespaces]
The \emph{formal support envelope} is the enumerated symbolic domain.
A \emph{canonical-parent carrier} requires a sourced inverse
aggregation.  The \emph{actual active support} requires an exact
nonzero coefficient certificate.  We treat these as three different
typed sets.
\end{definition}

Every current nonsoft row lies in the first namespace.  The other two
are not inferred from a symbolic coefficient expression.

\section{The cut-occurrence shadow}

\begin{definition}[Partial occurrence namespace]
For each \(c\in\Cns\), define
\[
 \iota(c):=(c,F_{\rm cut}(c))
\]
in a new typed namespace, and put
\[
 \Ocut:=\{\iota(c):c\in\Cns\}.
\]
The identifier of \(\iota(c)\) is derived deterministically from the
literal cut-path identifier.  It is called a \emph{partial
occurrence}; it is not an actual downstream occurrence.
\end{definition}

\begin{definition}[Canonical shadow matrix]
The cut-occurrence shadow is the linear map
\begin{equation}
 S_X:\mathbb C^{\Cns}\longrightarrow\mathbb C^{\Ocut},
 \qquad S_Xe_c=e_{\iota(c)}.
 \label{eq:shadow}
\end{equation}
Its only nonzero coefficient in column \(c\) is the exact rational
number \(1\).
\end{definition}

\begin{theorem}[Exact shadow identities]
\label{thm:shadow}
The matrix \(S_X\) is total on \(\Cns\), injective and row separated.
Moreover,
\begin{equation}
 \one_{\Ocut}^{T}S_X=\one_{\Cns}^{T}.
 \label{eq:shadow-conservation}
\end{equation}
On every nonzero edge it intertwines the native tuple, source
\(h_0\), physical normalization, source path and formal-support
status.
\end{theorem}

\begin{proof}
The rule \(c\mapsto\iota(c)\) creates one distinct typed row for each
literal source identifier, hence is injective and total.  Each column
contains one coefficient equal to \(1\), proving
\eqref{eq:shadow-conservation}.  The metadata statement follows
because \(F_{\rm cut}(c)\) is copied without reinterpretation from
the corresponding source obligation.
\end{proof}

\begin{remark}[Why the identity is useful but not the goal]
An identity matrix can be trivial algebraically and still useful as a
typed boundary object.  Here it gives a lossless handoff point: any
later construction must explain how one partial row branches into
actual rows without changing the sourced column total.  It does not
explain that branching.
\end{remark}

\subsection{Exact reconnection with the soft rows}

Let \(\mathcal N_X\) be the native-record set and write the
conservative cut matrix as
\[
 M_X^{\rm cut}=
 \begin{pmatrix}
 M_X^{\EPS}\\ M_X^{\rm ns}
 \end{pmatrix}.
\]

\begin{proposition}[Shadow-level reconnection]
\label{prop:reconnection}
The map \(I_{\EPS}\oplus S_X\) satisfies
\begin{equation}
 \one_{\EPS}^{T}M_X^{\EPS}
 +\one_{\Ocut}^{T}S_XM_X^{\rm ns}
 =\one_{\mathcal N_X}^{T}.
 \label{eq:reconnection}
\end{equation}
This is an exact coefficient identity at the cut-shadow level.
\end{proposition}

\begin{proof}
Insert \eqref{eq:shadow-conservation} into the cut-column identity
\(\one_{\EPS}^{T}M_X^{\EPS}
+\one_{\Cns}^{T}M_X^{\rm ns}=\one_{\mathcal N_X}^{T}\).
\end{proof}

\Cref{prop:reconnection} neither supplies physical cover nor proves
that a symbolic column is nonzero.  It only confirms that no
cut-level coefficient was deleted by introducing the shadow.

\section{Universal pushforward property}

The word ``maximal'' below refers to sourced cut fields, not to an
unknown actual carrier.

\begin{definition}[Actual extension candidate]
An actual extension candidate consists of a typed row set \(\Oact\),
a row-separated matrix
\[
 L_X:\mathbb C^{\Cns}\longrightarrow\mathbb C^{\Oact},
\]
and a forgetting map \(q:\Oact\to\Ocut\).  The map \(q\) must erase
only downstream fields and must retain all fields in
\(F_{\rm cut}\).  Its linear pushforward is
\[
 q_*e_o=e_{q(o)}.
\]
The candidate lies over the cut shadow when, for all \(c,c'\in\Cns\),
\begin{equation}
 \sum_{\substack{o\in\Oact\\q(o)=\iota(c)}}(L_X)_{o,c'}
 =\delta_{c,c'}.
 \label{eq:fiber-conservation}
\end{equation}
\end{definition}

\begin{theorem}[Conditional universal property]
\label{thm:universal}
Every actual extension candidate lying over the cut shadow satisfies
\begin{equation}
 q_*L_X=S_X.
 \label{eq:universal}
\end{equation}
Conversely, \eqref{eq:universal} is equivalent to the fiber identities
\eqref{eq:fiber-conservation}.
\end{theorem}

\begin{proof}
The coefficient of \(e_{\iota(c)}\) in \(q_*L_Xe_{c'}\) is the
left-hand side of \eqref{eq:fiber-conservation}.  It equals
\(\delta_{c,c'}\), which is exactly the coefficient of
\(e_{\iota(c)}\) in \(S_Xe_{c'}\).  Reading the same argument
backwards proves the converse.
\end{proof}

\begin{corollary}[Conservation inherited from an actual lift]
If \(q_*L_X=S_X\), then
\[
 \one_{\Ocut}^{T}q_*L_X=\one_{\Cns}^{T}.
\]
If, in addition, forgetting preserves the coefficient sum
\(\one_{\Ocut}^{T}q_*=\one_{\Oact}^{T}\), then
\(\one_{\Oact}^{T}L_X=\one_{\Cns}^{T}\).
\end{corollary}

The existence of \(q,L_X\), the number of rows in each fiber, and all
actual edge coefficients remain hypotheses.  Therefore
\cref{thm:universal} is a compatibility condition for future work,
not a construction of the missing lift.

\section{What remains deliberately absent}

The partial row reserves, but does not populate, the following field
groups:
\begin{enumerate}[leftmargin=2.2em]
\item canonical parent, determinant fiber and \(Q_D\) edge;
\item ordered outer zero mode and \(Q_Z\) edge;
\item physical occurrence, physical group, cover and reconnection;
\item transformation-stage identifier, exact later multiplier,
source and target shift tags, and downstream selector edge; and
\item the affine-corridor consumer fields and endpoint-ledger token.
\end{enumerate}
The TPC-141 occurrence registry names administrative stages, but it
does not give a row-level cut-to-stage-to-parent occurrence crosswalk
\citep{WangTPC141}.  It is therefore not used to manufacture these
fields.

\begin{proposition}[No promotion by shadow construction]
\label{prop:no-promotion}
Constructing \(S_X\) proves none of the following:
\begin{enumerate}[label=(\alph*),leftmargin=2em]
\item existence or uniqueness of an actual occurrence lift;
\item actual one-to-many branch multiplicities;
\item determinant, zero-mode, physical-group or downstream-shift
totality;
\item membership in actual active support; or
\item an arithmetic estimate for the nonsoft scalar.
\end{enumerate}
\end{proposition}

\begin{proof}
All five statements require fields erased by the map from a
hypothetical actual extension to \(\Ocut\).  The definition of
\(S_X\) uses only \(F_{\rm cut}\), so none is a consequence of
\eqref{eq:shadow}.
\end{proof}

\section{Executable finite certificate}

The standard-library program
\begin{center}
\small
\texttt{experiments/tpc153\_cut\_occurrence\_shadow.py}
\end{center}
reruns the TPC-143 generator, compares its logical UTF-8/LF
obligations and certificate, and emits one shadow record per source
obligation.  The lock mode is
\texttt{CANONICAL\_UTF8\_LF\_V2}; hashes have integrity semantics
only.

For the committed production fixture,
\[
 \#\Cns=2988,\qquad
 \#\mathcal C_X^{\ETO}=0,\qquad
 \#\mathcal C_X^{\FUM}=2988.
\]
All \(2988\) source columns have one shadow row and exact column sum
\(1\).  The validator rejects deletion, duplication, terminal
relabeling, nonunit weight, \(h_0\) or normalization drift, active
support promotion, an invented downstream selector, an invented
actual completion, and a false \(\Ltwo\) claim.

The production sample cannot exercise the \(\ETO\) branch.  A
separate regression uses the TPC-135 block policy at
\[
 X=2^{84},\quad R=2^{21},\quad V=2^{10},\quad
 (j_L,j_K)=(38,46),\quad D_0=2,
\]
for which the policy classifies the block as eligible.  This record is
marked \texttt{SYNTHETIC\_L0\_ONLY}.  It tests that the union
\eqref{eq:nonsoft-domain} remains typed, but proves neither a current
production \(\ETO\) row nor asymptotic \(\ETO\) existence.

\begin{center}
\small
\begin{tabularx}{\textwidth}{@{}P{0.38\textwidth}P{0.23\textwidth}Y@{}}
\toprule
Export & Status & Exact meaning\\
\midrule
\(\mathsf{H1.cut\_occurrence\_shadow}\)
& \(\PROVED_{\Lone}\) & Exact total conservative partial lift.\\
\(\mathsf{H1.frontier\_occurrence\_lift}\)
& \(\NT\) & Actual downstream rows and edge coefficients absent.\\
Current-schema-only actual-lift derivation
& \(\STOP\) & A shadow cannot supply erased provenance.\\
Selected augmented occurrence route
& Open & A theorem-backed external crosswalk may extend the shadow.\\
\bottomrule
\end{tabularx}
\end{center}

\section{Kill criteria and next artifact}

The shadow subroute fails immediately if any nonsoft source lacks one
partial row, if two source rows share a partial identifier, if a
column sum differs from \(1\), or if native, \(h_0\), normalization
or support metadata drift.  Conversely, passing these checks must
not be promoted to the actual lift.

The next required artifact is a theorem-backed, row-level
cut-to-stage-to-parent-to-physical-occurrence crosswalk on every
\(\ETO\) and \(\FUM\) source.  It must include exact edge
multipliers, native and \(h_0\) lineage, physical normalization,
support status, cover and reconnection.  The source-locked route
decision remains \(\NT\) until that object or the separately typed
scalar-plus-\(\ETO\) alternative is supplied \citep{WangTPC152}.

This paper proves no frontier \(o(X)\) estimate, positive
fixed-\(h_0\) \(\Ltwo\) saving, endpoint below \(1/400\), prime-pair
lower bound or twin-prime theorem.

\begingroup
\small
\raggedright
\bibliographystyle{plainnat}
\bibliography{references}
\endgroup
\end{document}
